% Ad Familiares 13.18 --- headnote
% ad-familiares-13-18, 46 BC, Rome

\textit{Cicero to Servius Sulpicius Rufus, proconsul of
Achaia, written from Rome around the beginning of November
46 BC (the manuscript dateline: \textit{Scr. Romae, ut
videtur, a. 708 (46)}). Part of the running Servius
sequence (Fam.\ 13.17--28); but unlike the formulaic
recommendations of that sheaf, this letter is purely
about Atticus and his property interests in Epirus.
Servius has written to Atticus on his own initiative,
offering his help with the management of those affairs in
the province; Atticus has just shown the letter to
Cicero, and Cicero, visibly delighted, writes back at
once. The piece is therefore the rare specimen in the
sheaf that conducts no commendation business but
acknowledges and reinforces a favour already volunteered.}

\textit{The structural play of the letter is the
self-conscious pivot in the second paragraph. Cicero
begins by listing the things he will not do --- he will
not ask Servius to act with redoubled zeal, and he will
not thank him for what he has already done unprompted ---
and then announces, on the strength of their old
friendship, that he will do both anyway. The closing
formula collapses the favour to Atticus into a personal
favour to Cicero himself: whatever Servius binds Atticus
by, he binds Cicero by as well. The first paragraph also
preserves a textual crux (the daggered phrase \textit{quod
tamen dubium nobis quin ita futurum fuerit non erat}),
where the manuscript reading is suspect; the rendering
here follows the sense common in the editors and treats
the parenthesis as a corroborative aside.}
